\documentclass[10pt,a4paper]{amsart}
\usepackage[margin=19mm]{geometry}
\usepackage{amsmath,amssymb,mathtools}
\usepackage[scr=rsfs,cal=euler]{mathalfa}
\usepackage{xfrac}
\usepackage[unicode,hidelinks,bookmarks=false]{hyperref}
\newcommand{\ie}{i.e.\ }
\newcommand{\cf}{cf.\ }
\newcommand{\resp}{respectively\ }
\newcommand{\Subsection}{\subsection}
\providecommand{\nobreakdash}{\nobreak\hbox{-}}
\setlength{\emergencystretch}{2em}
\multlinegap0pt
\usepackage{bm}
\usepackage{bbm}
\usepackage{dsfont}
\usepackage{xspace}
\usepackage{cancel}
\usepackage{mathtools}
\usepackage{amsthm}
\makeatletter
\@ifundefined{theorem}{\newtheorem{theorem}{Theorem}}{}
\@ifundefined{lemma}{\newtheorem{lemma}[theorem]{Lemma}}{}
\makeatother
\title[Asymptotics for multiple q-orthogonal polynomials from the R]{Asymptotics for multiple q-orthogonal polynomials from the RHP}
\author{Tomas Lasic Latimer}
\date{Source: arXiv:2502.00335v3 (CC BY 4.0)}
\begin{document}
\maketitle

\noindent\emph{Source and licence.} Asymptotics for multiple q-orthogonal polynomials from the RHP. Version arXiv:2502.00335v3 (CC BY 4.0), distributed under the Creative Commons Attribution 4.0 International licence, as recorded in the arXiv licence metadata for this exact version and shown on the abstract page https://arxiv.org/abs/2502.00335v3. Full rights notice: licenses/arxiv-2502.00335-CC-BY-4.0.txt.\par\smallskip

\noindent\textit{Editorial scope.} Counted unit: the Lemma whose source label is \texttt{None} in the article (source0.tex lines 490--492, source label \texttt{None}), delivered together with the article's own complete original proof of it (source0.tex lines 493--599, one \texttt{proof} environment of 107 lines). No mathematics is written, repaired or re-derived: statement and proof are the article's own text line for line. Required context retained uncounted: the main diff (lines 478--488). Every definition, display and label of those lines is retained unchanged. Omitted from this package and not redistributed: 1--477, 489--489, 600--1282. Nothing inside a retained excerpt is omitted or altered: every retained body is delivered exactly as the article's lines are written, so no omission receipt of any kind accompanies this package. Citations: the retained excerpts carry no citation command at all, so no bibliography entry of the article is transcribed and no reference is redistributed. Reference adaptation: none was needed and none was made; every reference inside the delivered unit resolves inside it, so no unresolved reference or citation is produced. Printed slips kept literal: no printed word, symbol, hypothesis or number of the counted statement or of its proof was altered. Layout and class: the article's own class is \texttt{amsart} with packages including geometry, biblatex, amsaddr, amsfonts, amsthm, amsmath, datetime, comment, enumitem, xcolor, tikz, tkz-graph, caption, tkz-euclide; the wrapper is the lane's pinned \texttt{amsart} editorial wrapper at 10pt on A4, which restates the theorem-like environments the retained text needs and adds the article's own macro definitions that the retained text needs (none), with extra wrapper packages (bm, bbm, dsfont, xspace, cancel, mathtools). The delivered numbering is the unit's own. Assets: no retained span contains a figure environment, a graphics-inclusion command, a PDF-page inclusion, a table or a TikZ picture; no image file is copied into this package, the package declares no asset dependency and no image file is redistributed with it. Licence: version arXiv:2502.00335v3 of this article is distributed under the Creative Commons Attribution 4.0 International licence, as recorded in the arXiv licence metadata for this exact version and shown on the abstract page https://arxiv.org/abs/2502.00335v3. A full rights notice is delivered at licenses/arxiv-2502.00335-CC-BY-4.0.txt. Characters: every retained excerpt is pure ASCII, so the delivered engine, XeLaTeX with its default Latin Modern text font, reports no missing character.
\begin{equation}\label{the main diff}
    Y_n(q^{-1}z) \begin{bmatrix}
   (1-z) &
   0 &
   0 \\
   0 &
   q^{\alpha} &
   0\\
   0&0&q^{\beta}
\end{bmatrix} = D_n(z)Y_n(z).
\end{equation}

\begin{lemma}
    The eigenvalues of $D_n(0)$ are given by $\lambda = 1,q^{\alpha},q^{\beta}$.
\end{lemma}

\begin{proof}
    The first column of Equation \eqref{the main diff} reads
    \[ Y_n^{(1)}(q^{-1}z)(1-z) = D_n(z)Y_n^{(1)}(z) .\]
    Let $U_n(z) = Y_n^{(1)}(z)(qz;q)_\infty$. It follows that 
    \[ U_n(q^{-1}z) = Y_n^{(1)}(q^{-1}z)(1-z)(qz;q)_\infty,\]
    and subsequently $U_n(z)$ satisfies the $q$-difference equation
    \[ U_n(q^{-1}z) = D_n(z)U_n(z) .\]
    We know that as $Y_n(z)$ solves RHP I, it has entries which are analytic in a neighbourhood about $z=0$, these are the entries which compose $Y_n(z)$ restricted to $\mathcal{D}_-$ (in fact we know that these entries are entire). Thus, as $(qz;q)_\infty$ is entire, the entries of $U_n(z)$ are also entire. Hence, $U_n(z)$ can be written as the converging power series:
\begin{equation}
    U_n(z) =  \sum_{i=0}^\infty\begin{bmatrix}
   a_i z^i\\
   b_i z^i\\
    c_i z^i
\end{bmatrix}.
\end{equation}
Substituting this series representation into the $q$-difference equation satisfied by $U_n(z)$ we find,
\begin{eqnarray*}
   \sum_{i=0}^\infty\begin{bmatrix}
   a_i q^{-i}z^i\\
   b_i q^{-i}z^i\\
    c_i q^{-i}z^i
 \end{bmatrix}&=& \begin{bmatrix}
    \mu_{1}-zq^{-2n} & \mu_{2} &\mu_{3}\\
   \mu_{4}&q^{n+\alpha}&0 \\
   \mu_{5}&0&q^{n+\beta}
   \end{bmatrix}\sum_{i=0}^\infty\begin{bmatrix}
   a_i z^i\\
   b_i z^i\\
    c_i z^i
 \end{bmatrix},\\
 0
 &=& \begin{bmatrix}
    \mu_{1}-q^{-i} & \mu_{2} &\mu_{3}\\
   \mu_{4}&q^{n+\alpha}-q^{-i}&0 \\
   \mu_{5}&0&q^{n+\beta}-q^{-i}
   \end{bmatrix}\sum_{i=0}^\infty\left(\begin{bmatrix}
   a_i z^i\\
   b_i z^i\\
    c_i z^i
 \end{bmatrix}+ \begin{bmatrix}
   a_i q^{-2n}z^{i+1}\\
   0\\
    0
 \end{bmatrix}\right),\\
  0
 &=& \begin{bmatrix}
    \mu_{1}-1 & \mu_{2} &\mu_{3}\\
   \mu_{4}&q^{n+\alpha}-1&0 \\
   \mu_{5}&0&q^{n+\beta}-1
   \end{bmatrix} \begin{bmatrix}
   a_0 \\
   b_0 \\
    c_0 
 \end{bmatrix}+ \mathcal{O}(z).
\end{eqnarray*}
Thus, for the above equation to be consistent we require that 1 is an eigenvalue of $D_n(0)$. Let us now turn our attention to the second column of Equation \eqref{the main diff}. The $q$-difference for this column reads
\[ Y_n^{(2)}(q^{-1}z)q^\alpha = D_n(z)Y_n^{(2)}(z).\]
Repeating our earlier arguments, we conclude that $Y_n^{(2)}(z)$ has entries which are entire and can be written as the converging power series
\begin{equation}
    Y_n^{(2)}(z) =  \sum_{i=0}^\infty\begin{bmatrix}
   a_i z^i\\
   b_i z^i\\
    c_i z^i
\end{bmatrix},
\end{equation}
where clearly $a_i$, $b_i$ and $c_i$ are different from $U_n(z)$. Substituting this series representation into the $q$-difference equation satisfied by $Y_n^{(2)}(z)$ we find,
\begin{eqnarray}
   q^{\alpha}\sum_{i=0}^\infty\begin{bmatrix}
   a_i q^{-i}z^i\\
   b_i q^{-i}z^i\\
    c_i q^{-i}z^i
 \end{bmatrix}&=& \begin{bmatrix}
    \mu_{1}-zq^{-2n} & \mu_{2} &\mu_{3}\\
   \mu_{4}&q^{n+\alpha}&0 \\
   \mu_{5}&0&q^{n+\beta}
   \end{bmatrix}\sum_{i=0}^\infty\begin{bmatrix}
   a_i z^i\\
   b_i z^i\\
    c_i z^i
 \end{bmatrix},\nonumber\\
 0
 &=& \begin{bmatrix}
    \mu_{1}-q^{-i+\alpha} & \mu_{2} &\mu_{3}\\
   \mu_{4}&q^{n+\alpha}-q^{-i+\alpha}&0 \\
   \mu_{5}&0&q^{n+\beta}-q^{-i+\alpha}
   \end{bmatrix}\sum_{i=0}^\infty\left(\begin{bmatrix}
   a_i z^i\\
   b_i z^i\\
    c_i z^i
 \end{bmatrix}+ \begin{bmatrix}
   a_i q^{-2n}z^{i+1}\\
   0\\
    0
 \end{bmatrix}\right),\nonumber\\
  0
 &=& \begin{bmatrix}
    \mu_{1}-q^{\alpha} & \mu_{2} &\mu_{3}\\
   \mu_{4}&q^{n+\alpha}-q^{\alpha}&0 \\
   \mu_{5}&0&q^{n+\beta}-q^{\alpha}
   \end{bmatrix} \begin{bmatrix}
   a_0 \\
   b_0 \\
    c_0 
 \end{bmatrix}+ \mathcal{O}(z).\label{Equation X}
\end{eqnarray}
For the above equation to be consistent we now require that $q^\alpha$ is an eigenvalue of $D_n(0)$. Repeating the same arguments for the third column it follows that $q^{\beta}$ must also be an eigenvalue $D_n(0)$.
\end{proof}

\end{document}
